\section{Love in the Time of Autumn}

Materialists, adhering to the theory of the ``selfish gene'', must regard Love as an epiphenomenon whose real purpose is the propagation of the genes. It is true that in the animal and vegetable kingdoms, at least in regards to the higher organisms, sex serves the purpose of the continuation of the species. Although \textbf{Vladimir Solovyov}, in the Meaning of Love, does identify sexual love as the basis for all love, a great love cannot be reduced to its biological function. In respect to human life, sex, more than simply the propagation of the species, has the task of elevating and perfecting human nature. Specifically, while this task demands the greatest possible diversity of specimens of humanity, more importantly there should appear in the world the best examples of humanity, both for themselves and for their influence on others.

Love does not obviously have any relationship to that task. Great men are most often born from ordinary parents without great passion. And the most passionate loves, are often barren. Solovyov provides several literary examples, such as Romeo and Juliet, whose passion ended in death without offspring. Anyone can produce posterity with any person of the opposite sex. Love, on the other hand, is satisfied only with a specific one.

In the wake of the attack in Kenya, a lady on CNN asked the expert how we could prevent a similar terrorist attack on a mall in the USA. That's easy: enforce federal background checks, close the ``gun show loophole'', ban assault rifles and large clips. Duh. At least, that is what ``they'' have been telling us.

Ron Shaich, the CEO of Panera Bread, tried to limit his grocery shopping to what someone gets on food stamps, or EBT cards. Of course, he was shocked how difficult it was to eat within those limits. Certainly, he couldn't eat at his own restaurants with their over-priced sandwiches. I saw him interviewed on Al Jazeera America and kept waiting for the interlocutor to ask him what he paid his store employees and whether they could afford to live on that salary. The question never came, but I suspect the answer is ``no''.

Pope Francis wants to deemphasize the Church's obsession with abortion, gay marriage, and contraception. The result will most likely be that, as issues, they will suffer from benign neglect without actually overturning Church teaching. This is what happened with usury, which is still prohibited by the Church, yet is almost never spoken about. Keep in mind that usury is the taking of any interest.

Pope Benedict envisioned a smaller, purer Church. Pope Francis instead wants a more inclusive, larger Church, peopled by the ``poor and the marginalized''. The successful, creative, and knowledgeable will continue to leave the Church on the margins.

\emph{Who would willingly be an aphid in a colony of ants?}


\paragraph{Completing the Task}


It is one thing to claim the identity of all traditions, but it is another to demonstrate it. Likewise, one can agree with Guenon or Schuon that Christianity is part of Tradition, it is another to show it. As they say, ``The devil is in the details.'' An important task in our project is to actually demonstrate it by pointing out neglected metaphysical teachings or spiritual practices, and then comparing them to other traditions. I think we have gone some distance in that task, or at least indicated fruitful paths for further research. What others do with this knowledge is not up to us.

In particular, the Medieval tradition has deep roots in the paganism of Greece and Rome. We have shown this particularly in the metaphysical understanding of ultimate reality and in the understanding of the inner structure of the soul. These roots also reappear in the outward hierarchical structure of society. We have shown that there is indeed a living esoteric tradition in the West, even if not so formally structured as some would like.

We have also delved into the Vedanta, pointing out that the experiences of the saints in the West are similar to those in the East, after taking account of the wide cultural gulf and modes of expression. Following Coomaraswamy, we agree that the great wisdom and spiritual literature of the West, such as Dante, Eckhart, St John, and so on, are sufficient in themselves, so recourse to alien forms of thought and customs are really unnecessary for serious spiritual seekers.

Yet these explorations have accentuated some differences. Fundamentally, for a mode of thought that accepts hierarchy and rejects egalitarianism, must we then suspend that point of view and regard all traditions as equal and arbitrary? Perhaps not.

For example, we have pointed out that the understanding of the Trinity demands that being and the Logos are one and inseparable. This is not compatible with the pagan idea that both the gods and the world arose from chaos and are subject to impersonal fate. No, logos, not chaos, is fundamental and ultimate; God is not Himself subject to fate, but is perfectly free. The world is directed by personal Providence, not by the Fates. For the pagan, there is eternal return, the same thing over and over again, without purpose, indifferent to human ends. This is not so, because Providence pulls events in a specific direction. Even Guenon acknowledges this in the \emph{Symbolism of the Cross}.

In the Medieval tradition, every man and woman is assigned his allotted place by divine Justice. He is free to act on his given talents to the limits of his abilities and opportunities. Contrast that to the stifling belief in reincarnation. In that teaching, the circumstances of birth is not a gift, but rather an abstract and impersonal punishment for some alleged transgressions in a previous life. One simply accepts his fate. I know Guenon claims that reincarnation is not the true teachings of the Vedanta, but try to find some guru in the ``still living tradition'' who denies it. Unless you can do so consistently and universally, that teaching has been lost.

Julius Evola often wrote of the ``dignity of the human person'', and expression taken from Catholic social teaching, but I don't know that it can be traced back to any explicit pagan teachings. It is certainly more compatible with the ideal of the absolute self rather than the annihilation of the self emphasized so often by Guenon as the Supreme Identity. I believe that the true identity is the recognition of the common nature following theosis, not the identity of the Self in the Vedanta. That identity would seemingly require the identify not only of essence, but also of existence, something impossible. Guenon avoids that impossibility by claiming that the identity is ``beyond Being'', i.e., it is non-being. Yet it is described as Sat-Chit-Ananda, i.e., Being-Consciousness-Bliss.

\paragraph{True Science of Man}


Beyond all the trivial and risible critiques of Tradition, ultimately it rests on whether it provides the true science of man, truer than the various scientific, psychological, or even religious theories. This science says that man has a body and a soul with different levels. In part, it agrees with science. As a body, man is subject to the laws of physics. As a living being, man is determined by biology, genetics, the laws of growth, etc. As an animal, man has kinship with the animals. Yet, man has a higher soul, or spirit, that transcends all those, and gives him access to intelligence, rationality, and the realms beyond the merely physical.

However, man has lost that last ability, he has lost ready access to the preternatural and supernatural aspects of his being, and appears, for the most part, as a natural being. This is what deceives the scientist and the materialist, who can only recognize the natural man.

Through spiritual and ascetical practice, man can seek to reintegrate the higher parts of his being. Of course, to refuse to engage in that practice is to be forever deaf and blind to the preternatural and supernatural aspects. That is the situation we are in.


\flrightit{Posted on 2013-09-24 by Cologero}

\begin{center}* * *\end{center}

\begin{footnotesize}
\begin{sffamily}

\texttt{Santiago on 2013-09-24 at 15:08 said:}

Thank you for the above post. Hopefully the following is helpful, I offer it in a constructive spirit:

If in discussing the identity of the Supreme with a being's existence (as well as, but distinct from, its essence) the title post refers to the view that a particular being's discrete conditions -- its empirical ego and body -- are God, and that therefore in realising this identity particularity is extinguished (``annihilation of the self''), then this opinion should perhaps not be attributed to the Vedanta. Atman is Brahman, and Brahman is Sat-Chit-Ananda. But ``Sat'' is Being in the absolute, not a particular individual ``jiva''.

Although, in a certain (holographic) sense, ``Jiva'' is ``Atman'', if the former is the ``drop'' in Rumi's formulation: ``You are not a drop in the ocean. You are the entire ocean in a drop''. The Bhagavad Gita refers to this: ``Verily the embodied living entity is My infinitesimal potency and eternal'' (15:7). Yet this does not mean Arjuna ceases to have his own particular body and dharma to fulfil. Precisely identity with God and the eternity of the ``embodied living entity'' are both true -- just as Christ is both one with His Father and at His Father's right hand.

In the same way, Christianity distinguishes between the Holy Trinity and creation, and yet Meister Eckhart identities the Trinity with creation (Expositio Libri Genesis) and particular beings with the Word, Christ (Sermon 69). What is considered illusory and extinguished (similarly to Vedanta) is a thing (mis)perceived as distinct from the Word. This is consistent with St. Clement of Alexandria: ``the Son is neither simply one thing, nor many things as parts, but one thing as all things, whence also he is all things'' (Stromata, IV.25).

\hfill

\texttt{Scardanelli on 2013-09-24 at 22:06 said:}

``BTW, Francis is the first pope without valid orders, a layman who can only emphasize socialist agenda. Sadly the Church as represented by Rome is done.''

Perhaps the task is not to pronounce the death knell of the Church of Rome, as so many are fond of doing, but to provide the enlivening impulse that revivifies her. Tomberg elaborates on this in Lazarus, Come Forth.

\hfill

\texttt{Synodius on 2013-09-25 at 07:23 said:}

I meant current representation in Vatican, not the whole Catholic tradition. Such an impulse can only come from the traditionalist camp (or eastern uniates if they realise what's going on). Not all of them are preaching `no salvation outside our temple', after all Rama Coomaraswamy was ordained in a sedevacantist lineage and his perennialism was no secret.

\hfill

\end{sffamily}
\end{footnotesize}

